\originalchapter{8.512作业4--8：电导、无序与超导}
\psource{Patrick A. Lee，8.512 Spring 2009，按网页作业编号；PDF04内部底稿写作业5、PDF05内部底稿写作业4，修订标记恢复网页编号。}
\originalsection{作业4（修订为2009年3月9日）}
\begin{exercise}[原题1(a)--(c)]\label{p:512-4-1}
电子气受$E(r,t)=E_0\e^{\ii(qr-\omega t)}$，弛豫时间Boltzmann方程$\partial_tf+v\cdot\nabla_rf-eE\cdot\nabla_pf=-(f-f_0)/\tau$。(a)取$f=f_0+\Phi(k)\e^{\ii(qr-\omega t)}$，一阶求沿$\hat e\parallel E_0$的电导。(b)由光子吸收与发射的能量耗散导出Kubo谱表达，并用Kramers--Kronig恢复复电导。(c)假设电子态是平面波，$|q|\ll k_F$且展宽$\eta=\hbar/\tau$，说明谱表达如何化为(a)。
\end{exercise}
\begin{solution}
(a)一阶式$[\tau^{-1}-\ii\omega+\ii q\cdot v]\Phi=e(E_0\cdot v)f_0'$，再用$j=-2e\int v\Phi\dd^3k/(2\pi)^3$，得
\begin{equation}
\sigma(q,\omega)=2e^2\int\frac{\dd^3k}{(2\pi)^3}\frac{\tau(\hat e\cdot v_k)^2[-f_0'(\epsilon_k)]}{1-\ii\tau(\omega-q\cdot v_k)}.
\end{equation}
(b)以实场$E_0\cos\omega t$，$A$幅度$E_0/\omega$，跃迁率的正频项由Golden rule给出。吸收减发射使权重$f_a-f_b$，平均耗散$\sigma'E_0^2/2$给
\begin{equation}
\sigma'(q,\omega)=\frac{2\pi e^2}{m^2V\omega}\sum_{ab}(f_a-f_b)|\avg{b|\e^{\ii qr}\hat e\cdot p|a}|^2\delta(E_b-E_a-\hbar\omega).
\end{equation}
这里2为自旋，$\pi\delta$采用通常归一。将$\pi\delta(x)$视为$\operatorname{Im}[-(x+\ii0)^{-1}]$，恢复满足因果性且含合适零频项的复响应；顺、抗磁两部分的结合见第\ref{ch:disorder}章。若用原题将$\pi$吸入谱$\delta$的记号，两式必须同时换约定。

(c)矩阵元要求$k_b=k_a+q$，小$q$下$\Delta E=\hbar q\cdot v$、$f_k-f_{k+q}=-\hbar q\cdot v f_0'$。将谱极点展宽、考虑$\omega=0$所需的静态抵消项后，得到$\tau/[1-\ii\tau(\omega-q\cdot v)]$。严格任意纵向$q$的守恒碰撞核还需局域化学势修正；原题的单$\tau$模型在此按给定近似使用。
\end{solution}
\begin{exercise}[原题2(a)--(b)]\label{p:512-4-2}
各向同性低温金属中取$E\perp q$。(a)从上式证明$\sigma_\perp=\sigma_0(1-\ii\omega\tau)^{-1}(3/4)\int_{-1}^1(1-x^2)/(1+sx)\dd x$，$s=\ii qv_F\tau/(1-\ii\omega\tau)$，$\sigma_0=ne^2\tau/m$。(b)求小$|s|$和大$|s|$的极限，并用粒子运动解释低频大$q\ell$时$1/(q\ell)$压低因子。
\end{exercise}
\begin{solution}
(a)令$q\parallel z,E\parallel x$，$v_x^2=v_F^2(1-x^2)\cos^2\phi$，对方位角积分$\pi$，径向积分由$-f'=\delta(E-E_F)$固定费米面。用$n/m=D(E_F)v_F^2/3$，得到所给归一。

(b)小$s$展开$1/(1+sx)$，奇次积分为0，$\int(1-x^2)\dd x=4/3$、$\int(x^2-x^4)\dd x=4/15$，故$\sigma_\perp=\sigma_0[1+s^2/5+O(s^4)]/(1-\ii\omega\tau)$。低频$s=\ii Q,Q=q\ell\gg1$时，实积分$\int(1-x^2)/(1+Q^2x^2)\dd x\simeq\pi/Q$，所以$\sigma_\perp\simeq3\pi\sigma_0/(4q\ell)$。只有$|v_z|/v_F\lesssim1/(q\ell)$的电子在一段平均自由程内不跨越许多反向电场周期，其速度主要平行电场；这一薄角带占总角空间$O(1/q\ell)$，给相同压低尺度。任意复杂$s$的大值展开应保留解析支，不能直接将这个实低频极限推广到所有频率。
\end{solution}
\originalsection{作业5（网页编号5，底稿题页写2008年3月10日）}
\begin{exercise}[原题1(a)--(f)，(f)选做]\label{p:512-5-1}
一维紧束缚跃迁$t$，随机站势$V_i\sim U[-W,W]$。原题把站$N$与站1认作同一站，所以独立站数$M=N-1$。(a)写周期$M\times M$矩阵。(b)改为$\psi_N=\e^{\ii\phi}\psi_1$。(c)证明与每键跃迁$t\e^{\ii\phi/M}$的周期模型等价，求穿环磁通。(d)取$W/t=2,N=20$，数值画零能附近10级随$\phi$变化，增大$N$比较。(e)求Thouless量$g=E_T/\delta$，$E_T=\partial_\phi^2E_\alpha|_0$，按附近10级与多随机样本平均，检查$N$依赖。(f)选做二维方格，比较$W/t=4,9$。
\end{exercise}
\begin{solution}
(a)$H_{ii}=V_i$、$H_{i,i+1}=H_{i+1,i}=t$，角元$H_{M,1}=H_{1,M}=t$。(b)取$H_{M,1}=t\e^{\ii\phi},H_{1,M}=t\e^{-\ii\phi}$，总绕环相位$\phi$。(c)局域变换$\psi_j=\e^{\ii j\phi/M}\widetilde\psi_j$把相位均匀分配到每键，净相位不变，$\Phi/\Phi_0=\phi/(2\pi)$，此处电子通量量子$\Phi_0=hc/e$。

(d)--(e)随卷脚本固定种子851205，直接对Hermitian矩阵本征值求解，取$N=20,40,80$等尺寸及多个独立样本。按能量选零附近10级，曲率用$[E(\delta\phi)-2E(0)+E(-\delta\phi)]/\delta\phi^2$，另用更小步长检查稳定性。能级曲率有正负；若直接有符号平均可因相消趋零，通常的正电导尺度须平均$|E''|$或RMS。证据同时保留有符号值与$|E''|$，$\delta$由所选能级跨度除9估计，$g=\avg{|E''|}/\avg\delta$。图和数据为该有限样本实验，局域化趋势由$g$随$L$衰减判断，单一$N=20$不能精确给无限系统指数。

(f)二维采用$L_x\times L_y$环面，一方向穿相位，另一方向周期；每站4邻居且每键Hermitian配对。源码计算多个小尺寸在$W/t=4,9$的曲率分布与$g$，强无序边界敏感性更弱。有限尺寸对照支持方法复现，不冒称数值证明二维无穷尺寸局域化。
\end{solution}
\begin{center}\begin{tikzpicture}\begin{axis}[width=11cm,height=5.7cm,xlabel={$\phi$},ylabel={$E/t$},title={一维无序环：$N=20$，10个能级},]
\addplot[thick,color=structurecolor] coordinates {(-3.1415927,-1.2419138) (-3.0630528,-1.2419166) (-2.984513,-1.2419247) (-2.9059732,-1.2419381) (-2.8274334,-1.2419568) (-2.7488936,-1.2419807) (-2.6703538,-1.2420095) (-2.5918139,-1.2420432) (-2.5132741,-1.2420815) (-2.4347343,-1.2421241) (-2.3561945,-1.2421708) (-2.2776547,-1.2422213) (-2.1991149,-1.2422753) (-2.120575,-1.2423324) (-2.0420352,-1.2423923) (-1.9634954,-1.2424547) (-1.8849556,-1.242519) (-1.8064158,-1.242585) (-1.727876,-1.2426522) (-1.6493361,-1.2427202) (-1.5707963,-1.2427885) (-1.4922565,-1.2428569) (-1.4137167,-1.2429247) (-1.3351769,-1.2429917) (-1.2566371,-1.2430574) (-1.1780972,-1.2431214) (-1.0995574,-1.2431833) (-1.0210176,-1.2432427) (-0.9424778,-1.2432993) (-0.86393798,-1.2433528) (-0.78539816,-1.2434027) (-0.70685835,-1.2434489) (-0.62831853,-1.243491) (-0.54977871,-1.2435287) (-0.4712389,-1.2435619) (-0.39269908,-1.2435903) (-0.31415927,-1.2436138) (-0.23561945,-1.2436323) (-0.15707963,-1.2436455) (-0.078539816,-1.2436535) (0,-1.2436561) (0.078539816,-1.2436535) (0.15707963,-1.2436455) (0.23561945,-1.2436323) (0.31415927,-1.2436138) (0.39269908,-1.2435903) (0.4712389,-1.2435619) (0.54977871,-1.2435287) (0.62831853,-1.243491) (0.70685835,-1.2434489) (0.78539816,-1.2434027) (0.86393798,-1.2433528) (0.9424778,-1.2432993) (1.0210176,-1.2432427) (1.0995574,-1.2431833) (1.1780972,-1.2431214) (1.2566371,-1.2430574) (1.3351769,-1.2429917) (1.4137167,-1.2429247) (1.4922565,-1.2428569) (1.5707963,-1.2427885) (1.6493361,-1.2427202) (1.727876,-1.2426522) (1.8064158,-1.242585) (1.8849556,-1.242519) (1.9634954,-1.2424547) (2.0420352,-1.2423923) (2.120575,-1.2423324) (2.1991149,-1.2422753) (2.2776547,-1.2422213) (2.3561945,-1.2421708) (2.4347343,-1.2421241) (2.5132741,-1.2420815) (2.5918139,-1.2420432) (2.6703538,-1.2420095) (2.7488936,-1.2419807) (2.8274334,-1.2419568) (2.9059732,-1.2419381) (2.984513,-1.2419247) (3.0630528,-1.2419166) (3.1415927,-1.2419138)};
\addplot[thick,color=main] coordinates {(-3.1415927,-0.90317756) (-3.0630528,-0.90316488) (-2.984513,-0.90312691) (-2.9059732,-0.90306385) (-2.8274334,-0.90297604) (-2.7488936,-0.90286394) (-2.6703538,-0.90272814) (-2.5918139,-0.90256939) (-2.5132741,-0.90238852) (-2.4347343,-0.90218652) (-2.3561945,-0.9019645) (-2.2776547,-0.90172365) (-2.1991149,-0.90146532) (-2.120575,-0.90119095) (-2.0420352,-0.90090206) (-1.9634954,-0.90060031) (-1.8849556,-0.90028742) (-1.8064158,-0.8999652) (-1.727876,-0.89963555) (-1.6493361,-0.8993004) (-1.5707963,-0.89896177) (-1.4922565,-0.89862172) (-1.4137167,-0.89828233) (-1.3351769,-0.89794572) (-1.2566371,-0.89761401) (-1.1780972,-0.89728931) (-1.0995574,-0.89697373) (-1.0210176,-0.89666933) (-0.9424778,-0.89637813) (-0.86393798,-0.89610208) (-0.78539816,-0.89584306) (-0.70685835,-0.89560285) (-0.62831853,-0.89538309) (-0.54977871,-0.89518535) (-0.4712389,-0.89501101) (-0.39269908,-0.8948613) (-0.31415927,-0.89473732) (-0.23561945,-0.89463994) (-0.15707963,-0.89456987) (-0.078539816,-0.89452762) (0,-0.89451351) (0.078539816,-0.89452762) (0.15707963,-0.89456987) (0.23561945,-0.89463994) (0.31415927,-0.89473732) (0.39269908,-0.8948613) (0.4712389,-0.89501101) (0.54977871,-0.89518535) (0.62831853,-0.89538309) (0.70685835,-0.89560285) (0.78539816,-0.89584306) (0.86393798,-0.89610208) (0.9424778,-0.89637813) (1.0210176,-0.89666933) (1.0995574,-0.89697373) (1.1780972,-0.89728931) (1.2566371,-0.89761401) (1.3351769,-0.89794572) (1.4137167,-0.89828233) (1.4922565,-0.89862172) (1.5707963,-0.89896177) (1.6493361,-0.8993004) (1.727876,-0.89963555) (1.8064158,-0.8999652) (1.8849556,-0.90028742) (1.9634954,-0.90060031) (2.0420352,-0.90090206) (2.120575,-0.90119095) (2.1991149,-0.90146532) (2.2776547,-0.90172365) (2.3561945,-0.9019645) (2.4347343,-0.90218652) (2.5132741,-0.90238852) (2.5918139,-0.90256939) (2.6703538,-0.90272814) (2.7488936,-0.90286394) (2.8274334,-0.90297604) (2.9059732,-0.90306385) (2.984513,-0.90312691) (3.0630528,-0.90316488) (3.1415927,-0.90317756)};
\addplot[thick,color=blue] coordinates {(-3.1415927,-0.7379716) (-3.0630528,-0.73798646) (-2.984513,-0.73803096) (-2.9059732,-0.73810485) (-2.8274334,-0.73820773) (-2.7488936,-0.73833903) (-2.6703538,-0.73849804) (-2.5918139,-0.73868387) (-2.5132741,-0.7388955) (-2.4347343,-0.73913175) (-2.3561945,-0.7393913) (-2.2776547,-0.7396727) (-2.1991149,-0.73997436) (-2.120575,-0.74029456) (-2.0420352,-0.74063147) (-1.9634954,-0.74098315) (-1.8849556,-0.74134756) (-1.8064158,-0.74172255) (-1.727876,-0.74210593) (-1.6493361,-0.7424954) (-1.5707963,-0.7428886) (-1.4922565,-0.74328316) (-1.4137167,-0.74367665) (-1.3351769,-0.74406662) (-1.2566371,-0.74445063) (-1.1780972,-0.74482624) (-1.0995574,-0.74519104) (-1.0210176,-0.74554268) (-0.9424778,-0.74587884) (-0.86393798,-0.7461973) (-0.78539816,-0.74649595) (-0.70685835,-0.74677277) (-0.62831853,-0.74702586) (-0.54977871,-0.74725351) (-0.4712389,-0.74745414) (-0.39269908,-0.74762634) (-0.31415927,-0.74776893) (-0.23561945,-0.74788088) (-0.15707963,-0.74796143) (-0.078539816,-0.74800998) (0,-0.7480262) (0.078539816,-0.74800998) (0.15707963,-0.74796143) (0.23561945,-0.74788088) (0.31415927,-0.74776893) (0.39269908,-0.74762634) (0.4712389,-0.74745414) (0.54977871,-0.74725351) (0.62831853,-0.74702586) (0.70685835,-0.74677277) (0.78539816,-0.74649595) (0.86393798,-0.7461973) (0.9424778,-0.74587884) (1.0210176,-0.74554268) (1.0995574,-0.74519104) (1.1780972,-0.74482624) (1.2566371,-0.74445063) (1.3351769,-0.74406662) (1.4137167,-0.74367665) (1.4922565,-0.74328316) (1.5707963,-0.7428886) (1.6493361,-0.7424954) (1.727876,-0.74210593) (1.8064158,-0.74172255) (1.8849556,-0.74134756) (1.9634954,-0.74098315) (2.0420352,-0.74063147) (2.120575,-0.74029456) (2.1991149,-0.73997436) (2.2776547,-0.7396727) (2.3561945,-0.7393913) (2.4347343,-0.73913175) (2.5132741,-0.7388955) (2.5918139,-0.73868387) (2.6703538,-0.73849804) (2.7488936,-0.73833903) (2.8274334,-0.73820773) (2.9059732,-0.73810485) (2.984513,-0.73803096) (3.0630528,-0.73798646) (3.1415927,-0.7379716)};
\addplot[thick,color=gray] coordinates {(-3.1415927,-0.16642575) (-3.0630528,-0.16640191) (-2.984513,-0.16633051) (-2.9059732,-0.1662119) (-2.8274334,-0.16604666) (-2.7488936,-0.16583563) (-2.6703538,-0.16557984) (-2.5918139,-0.16528058) (-2.5132741,-0.16493936) (-2.4347343,-0.16455791) (-2.3561945,-0.16413819) (-2.2776547,-0.16368236) (-2.1991149,-0.16319281) (-2.120575,-0.16267213) (-2.0420352,-0.16212311) (-1.9634954,-0.16154872) (-1.8849556,-0.16095213) (-1.8064158,-0.16033668) (-1.727876,-0.15970586) (-1.6493361,-0.15906332) (-1.5707963,-0.15841284) (-1.4922565,-0.1577583) (-1.4137167,-0.15710371) (-1.3351769,-0.15645313) (-1.2566371,-0.15581067) (-1.1780972,-0.15518049) (-1.0995574,-0.15456674) (-1.0210176,-0.15397353) (-0.9424778,-0.15340491) (-0.86393798,-0.15286484) (-0.78539816,-0.15235716) (-0.70685835,-0.1518855) (-0.62831853,-0.15145333) (-0.54977871,-0.15106386) (-0.4712389,-0.15072001) (-0.39269908,-0.15042442) (-0.31415927,-0.15017935) (-0.23561945,-0.14998672) (-0.15707963,-0.14984802) (-0.078539816,-0.14976437) (0,-0.14973641) (0.078539816,-0.14976437) (0.15707963,-0.14984802) (0.23561945,-0.14998672) (0.31415927,-0.15017935) (0.39269908,-0.15042442) (0.4712389,-0.15072001) (0.54977871,-0.15106386) (0.62831853,-0.15145333) (0.70685835,-0.1518855) (0.78539816,-0.15235716) (0.86393798,-0.15286484) (0.9424778,-0.15340491) (1.0210176,-0.15397353) (1.0995574,-0.15456674) (1.1780972,-0.15518049) (1.2566371,-0.15581067) (1.3351769,-0.15645313) (1.4137167,-0.15710371) (1.4922565,-0.1577583) (1.5707963,-0.15841284) (1.6493361,-0.15906332) (1.727876,-0.15970586) (1.8064158,-0.16033668) (1.8849556,-0.16095213) (1.9634954,-0.16154872) (2.0420352,-0.16212311) (2.120575,-0.16267213) (2.1991149,-0.16319281) (2.2776547,-0.16368236) (2.3561945,-0.16413819) (2.4347343,-0.16455791) (2.5132741,-0.16493936) (2.5918139,-0.16528058) (2.6703538,-0.16557984) (2.7488936,-0.16583563) (2.8274334,-0.16604666) (2.9059732,-0.1662119) (2.984513,-0.16633051) (3.0630528,-0.16640191) (3.1415927,-0.16642575)};
\addplot[thick,color=orange] coordinates {(-3.1415927,0.1822732) (-3.0630528,0.18148306) (-2.984513,0.1795113) (-2.9059732,0.17694124) (-2.8274334,0.17409312) (-2.7488936,0.17111178) (-2.6703538,0.16806639) (-2.5918139,0.16499386) (-2.5132741,0.16191624) (-2.4347343,0.15884816) (-2.3561945,0.15580027) (-2.2776547,0.15278097) (-2.1991149,0.14979732) (-2.120575,0.14685554) (-2.0420352,0.14396134) (-1.9634954,0.1411201) (-1.8849556,0.13833698) (-1.8064158,0.135617) (-1.727876,0.1329651) (-1.6493361,0.13038616) (-1.5707963,0.12788501) (-1.4922565,0.12546647) (-1.4137167,0.1231353) (-1.3351769,0.12089625) (-1.2566371,0.11875404) (-1.1780972,0.11671329) (-1.0995574,0.11477859) (-1.0210176,0.11295442) (-0.9424778,0.11124513) (-0.86393798,0.10965494) (-0.78539816,0.10818789) (-0.70685835,0.10684782) (-0.62831853,0.10563834) (-0.54977871,0.10456276) (-0.4712389,0.10362413) (-0.39269908,0.10282513) (-0.31415927,0.10216812) (-0.23561945,0.10165504) (-0.15707963,0.10128743) (-0.078539816,0.10106642) (0,0.10099267) (0.078539816,0.10106642) (0.15707963,0.10128743) (0.23561945,0.10165504) (0.31415927,0.10216812) (0.39269908,0.10282513) (0.4712389,0.10362413) (0.54977871,0.10456276) (0.62831853,0.10563834) (0.70685835,0.10684782) (0.78539816,0.10818789) (0.86393798,0.10965494) (0.9424778,0.11124513) (1.0210176,0.11295442) (1.0995574,0.11477859) (1.1780972,0.11671329) (1.2566371,0.11875404) (1.3351769,0.12089625) (1.4137167,0.1231353) (1.4922565,0.12546647) (1.5707963,0.12788501) (1.6493361,0.13038616) (1.727876,0.1329651) (1.8064158,0.135617) (1.8849556,0.13833698) (1.9634954,0.1411201) (2.0420352,0.14396134) (2.120575,0.14685554) (2.1991149,0.14979732) (2.2776547,0.15278097) (2.3561945,0.15580027) (2.4347343,0.15884816) (2.5132741,0.16191624) (2.5918139,0.16499386) (2.6703538,0.16806639) (2.7488936,0.17111178) (2.8274334,0.17409312) (2.9059732,0.17694124) (2.984513,0.1795113) (3.0630528,0.18148306) (3.1415927,0.1822732)};
\addplot[thick,color=teal] coordinates {(-3.1415927,0.19531149) (-3.0630528,0.19610381) (-2.984513,0.19808209) (-2.9059732,0.20066302) (-2.8274334,0.20352633) (-2.7488936,0.20652716) (-2.6703538,0.20959631) (-2.5918139,0.21269684) (-2.5132741,0.21580665) (-2.4347343,0.21891104) (-2.3561945,0.22199928) (-2.2776547,0.2250629) (-2.1991149,0.22809473) (-2.120575,0.23108841) (-2.0420352,0.23403809) (-1.9634954,0.23693823) (-1.8849556,0.23978345) (-1.8064158,0.24256851) (-1.727876,0.24528821) (-1.6493361,0.24793737) (-1.5707963,0.25051083) (-1.4922565,0.25300341) (-1.4137167,0.25540992) (-1.3351769,0.25772517) (-1.2566371,0.25994396) (-1.1780972,0.26206112) (-1.0995574,0.26407151) (-1.0210176,0.26597007) (-0.9424778,0.26775181) (-0.86393798,0.26941187) (-0.78539816,0.27094557) (-0.70685835,0.27234841) (-0.62831853,0.27361615) (-0.54977871,0.27474484) (-0.4712389,0.27573086) (-0.39269908,0.27657096) (-0.31415927,0.27726231) (-0.23561945,0.27780256) (-0.15707963,0.27818982) (-0.078539816,0.27842272) (0,0.27850045) (0.078539816,0.27842272) (0.15707963,0.27818982) (0.23561945,0.27780256) (0.31415927,0.27726231) (0.39269908,0.27657096) (0.4712389,0.27573086) (0.54977871,0.27474484) (0.62831853,0.27361615) (0.70685835,0.27234841) (0.78539816,0.27094557) (0.86393798,0.26941187) (0.9424778,0.26775181) (1.0210176,0.26597007) (1.0995574,0.26407151) (1.1780972,0.26206112) (1.2566371,0.25994396) (1.3351769,0.25772517) (1.4137167,0.25540992) (1.4922565,0.25300341) (1.5707963,0.25051083) (1.6493361,0.24793737) (1.727876,0.24528821) (1.8064158,0.24256851) (1.8849556,0.23978345) (1.9634954,0.23693823) (2.0420352,0.23403809) (2.120575,0.23108841) (2.1991149,0.22809473) (2.2776547,0.2250629) (2.3561945,0.22199928) (2.4347343,0.21891104) (2.5132741,0.21580665) (2.5918139,0.21269684) (2.6703538,0.20959631) (2.7488936,0.20652716) (2.8274334,0.20352633) (2.9059732,0.20066302) (2.984513,0.19808209) (3.0630528,0.19610381) (3.1415927,0.19531149)};
\addplot[thick,color=black] coordinates {(-3.1415927,0.51642451) (-3.0630528,0.51639826) (-2.984513,0.51631965) (-2.9059732,0.51618903) (-2.8274334,0.51600701) (-2.7488936,0.51577446) (-2.6703538,0.51549247) (-2.5918139,0.51516238) (-2.5132741,0.51478576) (-2.4347343,0.51436444) (-2.3561945,0.51390047) (-2.2776547,0.51339614) (-2.1991149,0.51285397) (-2.120575,0.51227671) (-2.0420352,0.51166732) (-1.9634954,0.51102899) (-1.8849556,0.51036513) (-1.8064158,0.50967932) (-1.727876,0.50897537) (-1.6493361,0.50825725) (-1.5707963,0.5075291) (-1.4922565,0.50679522) (-1.4137167,0.50606004) (-1.3351769,0.50532812) (-1.2566371,0.5046041) (-1.1780972,0.50389266) (-1.0995574,0.50319856) (-1.0210176,0.50252653) (-0.9424778,0.50188126) (-0.86393798,0.50126738) (-0.78539816,0.50068937) (-0.70685835,0.50015157) (-0.62831853,0.49965809) (-0.54977871,0.49921277) (-0.4712389,0.49881915) (-0.39269908,0.49848041) (-0.31415927,0.49819931) (-0.23561945,0.4979782) (-0.15707963,0.49781891) (-0.078539816,0.4977228) (0,0.49769067) (0.078539816,0.4977228) (0.15707963,0.49781891) (0.23561945,0.4979782) (0.31415927,0.49819931) (0.39269908,0.49848041) (0.4712389,0.49881915) (0.54977871,0.49921277) (0.62831853,0.49965809) (0.70685835,0.50015157) (0.78539816,0.50068937) (0.86393798,0.50126738) (0.9424778,0.50188126) (1.0210176,0.50252653) (1.0995574,0.50319856) (1.1780972,0.50389266) (1.2566371,0.5046041) (1.3351769,0.50532812) (1.4137167,0.50606004) (1.4922565,0.50679522) (1.5707963,0.5075291) (1.6493361,0.50825725) (1.727876,0.50897537) (1.8064158,0.50967932) (1.8849556,0.51036513) (1.9634954,0.51102899) (2.0420352,0.51166732) (2.120575,0.51227671) (2.1991149,0.51285397) (2.2776547,0.51339614) (2.3561945,0.51390047) (2.4347343,0.51436444) (2.5132741,0.51478576) (2.5918139,0.51516238) (2.6703538,0.51549247) (2.7488936,0.51577446) (2.8274334,0.51600701) (2.9059732,0.51618903) (2.984513,0.51631965) (3.0630528,0.51639826) (3.1415927,0.51642451)};
\addplot[thick,color=structurecolor] coordinates {(-3.1415927,0.92147543) (-3.0630528,0.9214817) (-2.984513,0.92150047) (-2.9059732,0.92153163) (-2.8274334,0.92157496) (-2.7488936,0.9216302) (-2.6703538,0.921697) (-2.5918139,0.92177493) (-2.5132741,0.92186348) (-2.4347343,0.92196212) (-2.3561945,0.92207019) (-2.2776547,0.92218704) (-2.1991149,0.9223119) (-2.120575,0.92244401) (-2.0420352,0.92258253) (-1.9634954,0.92272659) (-1.8849556,0.92287529) (-1.8064158,0.9230277) (-1.727876,0.92318288) (-1.6493361,0.92333986) (-1.5707963,0.92349767) (-1.4922565,0.92365534) (-1.4137167,0.92381189) (-1.3351769,0.92396638) (-1.2566371,0.92411784) (-1.1780972,0.92426536) (-1.0995574,0.92440804) (-1.0210176,0.92454501) (-0.9424778,0.92467543) (-0.86393798,0.92479854) (-0.78539816,0.92491357) (-0.70685835,0.92501983) (-0.62831853,0.9251167) (-0.54977871,0.92520358) (-0.4712389,0.92527996) (-0.39269908,0.92534537) (-0.31415927,0.92539943) (-0.23561945,0.92544182) (-0.15707963,0.92547228) (-0.078539816,0.92549063) (0,0.92549675) (0.078539816,0.92549063) (0.15707963,0.92547228) (0.23561945,0.92544182) (0.31415927,0.92539943) (0.39269908,0.92534537) (0.4712389,0.92527996) (0.54977871,0.92520358) (0.62831853,0.9251167) (0.70685835,0.92501983) (0.78539816,0.92491357) (0.86393798,0.92479854) (0.9424778,0.92467543) (1.0210176,0.92454501) (1.0995574,0.92440804) (1.1780972,0.92426536) (1.2566371,0.92411784) (1.3351769,0.92396638) (1.4137167,0.92381189) (1.4922565,0.92365534) (1.5707963,0.92349767) (1.6493361,0.92333986) (1.727876,0.92318288) (1.8064158,0.9230277) (1.8849556,0.92287529) (1.9634954,0.92272659) (2.0420352,0.92258253) (2.120575,0.92244401) (2.1991149,0.9223119) (2.2776547,0.92218704) (2.3561945,0.92207019) (2.4347343,0.92196212) (2.5132741,0.92186348) (2.5918139,0.92177493) (2.6703538,0.921697) (2.7488936,0.9216302) (2.8274334,0.92157496) (2.9059732,0.92153163) (2.984513,0.92150047) (3.0630528,0.9214817) (3.1415927,0.92147543)};
\addplot[thick,color=main] coordinates {(-3.1415927,1.3685779) (-3.0630528,1.3685752) (-2.984513,1.368567) (-2.9059732,1.3685535) (-2.8274334,1.3685347) (-2.7488936,1.3685107) (-2.6703538,1.3684817) (-2.5918139,1.3684478) (-2.5132741,1.3684093) (-2.4347343,1.3683664) (-2.3561945,1.3683193) (-2.2776547,1.3682684) (-2.1991149,1.3682139) (-2.120575,1.3681563) (-2.0420352,1.3680958) (-1.9634954,1.3680328) (-1.8849556,1.3679678) (-1.8064158,1.3679011) (-1.727876,1.367833) (-1.6493361,1.3677642) (-1.5707963,1.3676949) (-1.4922565,1.3676256) (-1.4137167,1.3675567) (-1.3351769,1.3674887) (-1.2566371,1.367422) (-1.1780972,1.3673569) (-1.0995574,1.3672939) (-1.0210176,1.3672334) (-0.9424778,1.3671757) (-0.86393798,1.3671213) (-0.78539816,1.3670703) (-0.70685835,1.3670233) (-0.62831853,1.3669803) (-0.54977871,1.3669418) (-0.4712389,1.3669079) (-0.39269908,1.3668789) (-0.31415927,1.3668548) (-0.23561945,1.366836) (-0.15707963,1.3668225) (-0.078539816,1.3668143) (0,1.3668116) (0.078539816,1.3668143) (0.15707963,1.3668225) (0.23561945,1.366836) (0.31415927,1.3668548) (0.39269908,1.3668789) (0.4712389,1.3669079) (0.54977871,1.3669418) (0.62831853,1.3669803) (0.70685835,1.3670233) (0.78539816,1.3670703) (0.86393798,1.3671213) (0.9424778,1.3671757) (1.0210176,1.3672334) (1.0995574,1.3672939) (1.1780972,1.3673569) (1.2566371,1.367422) (1.3351769,1.3674887) (1.4137167,1.3675567) (1.4922565,1.3676256) (1.5707963,1.3676949) (1.6493361,1.3677642) (1.727876,1.367833) (1.8064158,1.3679011) (1.8849556,1.3679678) (1.9634954,1.3680328) (2.0420352,1.3680958) (2.120575,1.3681563) (2.1991149,1.3682139) (2.2776547,1.3682684) (2.3561945,1.3683193) (2.4347343,1.3683664) (2.5132741,1.3684093) (2.5918139,1.3684478) (2.6703538,1.3684817) (2.7488936,1.3685107) (2.8274334,1.3685347) (2.9059732,1.3685535) (2.984513,1.368567) (3.0630528,1.3685752) (3.1415927,1.3685779)};
\addplot[thick,color=blue] coordinates {(-3.1415927,1.5897831) (-3.0630528,1.5897844) (-2.984513,1.5897883) (-2.9059732,1.5897948) (-2.8274334,1.5898039) (-2.7488936,1.5898155) (-2.6703538,1.5898294) (-2.5918139,1.5898458) (-2.5132741,1.5898643) (-2.4347343,1.589885) (-2.3561945,1.5899076) (-2.2776547,1.5899321) (-2.1991149,1.5899583) (-2.120575,1.589986) (-2.0420352,1.590015) (-1.9634954,1.5900453) (-1.8849556,1.5900765) (-1.8064158,1.5901085) (-1.727876,1.5901411) (-1.6493361,1.5901742) (-1.5707963,1.5902074) (-1.4922565,1.5902406) (-1.4137167,1.5902735) (-1.3351769,1.5903061) (-1.2566371,1.590338) (-1.1780972,1.5903691) (-1.0995574,1.5903992) (-1.0210176,1.5904281) (-0.9424778,1.5904557) (-0.86393798,1.5904817) (-0.78539816,1.590506) (-0.70685835,1.5905284) (-0.62831853,1.5905489) (-0.54977871,1.5905673) (-0.4712389,1.5905835) (-0.39269908,1.5905973) (-0.31415927,1.5906087) (-0.23561945,1.5906177) (-0.15707963,1.5906242) (-0.078539816,1.590628) (0,1.5906293) (0.078539816,1.590628) (0.15707963,1.5906242) (0.23561945,1.5906177) (0.31415927,1.5906087) (0.39269908,1.5905973) (0.4712389,1.5905835) (0.54977871,1.5905673) (0.62831853,1.5905489) (0.70685835,1.5905284) (0.78539816,1.590506) (0.86393798,1.5904817) (0.9424778,1.5904557) (1.0210176,1.5904281) (1.0995574,1.5903992) (1.1780972,1.5903691) (1.2566371,1.590338) (1.3351769,1.5903061) (1.4137167,1.5902735) (1.4922565,1.5902406) (1.5707963,1.5902074) (1.6493361,1.5901742) (1.727876,1.5901411) (1.8064158,1.5901085) (1.8849556,1.5900765) (1.9634954,1.5900453) (2.0420352,1.590015) (2.120575,1.589986) (2.1991149,1.5899583) (2.2776547,1.5899321) (2.3561945,1.5899076) (2.4347343,1.589885) (2.5132741,1.5898643) (2.5918139,1.5898458) (2.6703538,1.5898294) (2.7488936,1.5898155) (2.8274334,1.5898039) (2.9059732,1.5897948) (2.984513,1.5897883) (3.0630528,1.5897844) (3.1415927,1.5897831)};
\end{axis}\end{tikzpicture}\end{center}

\begin{center}\begin{tikzpicture}\begin{axis}[width=11cm,height=5.7cm,xlabel={$N$},ylabel={$g$},title={有限样本Thouless电导尺度},ymode=log,mark=*]
\addplot[thick,color=structurecolor] coordinates {(20,0.033114471) (40,0.0060952993) (80,7.5557181e-05) (140,2.5499047e-07)};
\addlegendentry{$W/t=2$}
\end{axis}\end{tikzpicture}\end{center}

\begin{center}\begin{tabular}{rrrr}\toprule $N$&独立站数&样本数&$g$\\\midrule
20&19&40&$0.03311$\\
40&39&40&$0.006095$\\
80&79&40&$7.56\times10^{-5}$\\\bottomrule\end{tabular}\end{center}
大尺寸的极小曲率逐渐接近浮点分辨率；图中更大尺寸只作带步长比较的数值边界记录，不用于拟合精确局域长度。
\begin{center}\begin{tikzpicture}\begin{axis}[width=11cm,height=5.7cm,xlabel={$L$},ylabel={$g$},title={二维有限尺寸对照},ymode=log,mark=*]
\addplot[thick,color=structurecolor] coordinates {(4,0.25939206) (6,0.13353791) (8,0.099962056)};
\addlegendentry{$W/t=4$}
\addplot[thick,color=main] coordinates {(4,0.023333335) (6,0.012444557) (8,0.0062368641)};
\addlegendentry{$W/t=9$}
\end{axis}\end{tikzpicture}\end{center}

\originalsection{作业6（底稿截止2008年3月31日）}
\begin{exercise}[原题1(a)--(b)]\label{p:512-6-1}
定义$J_\mu=-K_{\mu\nu}A_\nu$，将$K$分为横向与纵向投影。(a)由线性响应求自由费米气的静态$K_\perp(q,0)$，可取$q\parallel z$算$K_{xx}$。(b)证明双自旋Landau抗磁Gaussian式$\chi_D=-e^2k_F/(12\pi^2mc^2)$，是Pauli值$-1/3$。给原题$\hbar=1$之外的单位恢复。
\end{exercise}
\begin{solution}
(a)顺磁泡图与抗磁项相加，SI制可写
\begin{equation}
K_{xx}(q,0)=ne^2/m+(2e^2\hbar^2/m^2)\int\dd^3k\,k_x^2[f(\epsilon_{k+q/2})-f(\epsilon_{k-q/2})]/[(2\pi)^3(\epsilon_{k+q/2}-\epsilon_{k-q/2})]
\end{equation}
$q=0$项以积分分部$\int k_x\partial_{k_x}f=-\int f$恰抵消。为求$q^2$，令$b=\hbar^2q^2/(8m)$、$a=\hbar^2k_zq/(2m)$，比值为$f'(\epsilon+b)+a^2f'''(\epsilon+b)/6+\cdots$。把$b f''+a^2f'''/6$在球面作角积分，$\avg{k_x^2}=k^2/3$、$\avg{k_x^2k_z^2}=k^4/15$，径向的$\delta$导数积分给
\begin{equation}
K_\perp(q,0)=\frac{e^2k_F}{12\pi^2m}q^2+O(q^4)\qquad(\text{SI电流--}A\text{核}).
\end{equation}
(b)静态磁化电流$j=\nabla\times M$，线性$M=\chi_DB/\mu_0$，而$B=\nabla\times A$，横向Fourier给$j=(\chi_D/\mu_0)q^2A$，故$\chi_D=-\mu_0e^2k_F/(12\pi^2m)$。Pauli值$\chi_P=\mu_0\mu_B^2D(E_F)=\mu_0e^2k_F/(4\pi^2m)$，比值$-1/3$。Gaussian把$\mu_0$的磁化关系换成相应$c^{-2}$得到原式；裸质量与非相互作用假设不可省。
\end{solution}
\originalsection{作业7（2009年4月6日截止）}
\begin{exercise}[原题1(a)--(d)]\label{p:512-7-1}
无序超导体用正常态精确轨道$\phi_\alpha$表示$\int j_x^P=e\sum_{\alpha\beta\sigma}v_{\alpha\beta}c_{\beta\sigma}^\dagger c_{\alpha\sigma}$，$v_{\alpha\beta}=\int\phi_\beta^*p_x\phi_\alpha/m$，$\xi_\alpha=\epsilon_\alpha-\mu$，$E_\alpha=\sqrt{\xi_\alpha^2+\Delta^2}$，$u_\alpha^2=(1+\xi_\alpha/E_\alpha)/2$，$v_\alpha^2=(1-\xi_\alpha/E_\alpha)/2$，$\hbar=1$。(a)证明零温光吸收只创建两个准粒子，权重$(u_\alpha v_\beta-v_\alpha u_\beta)^2|v_{\alpha\beta}|^2\delta(\omega-E_\alpha-E_\beta)$，并以$f(\xi,\xi')=\Omega^{-1}\sum|v_{\alpha\beta}|^2\delta(\xi-\xi_\alpha)\delta(\xi'-\xi_\beta)$写积分。(b)证明相干因子$[1-\xi\xi'/(EE')-\Delta^2/(EE')]/2$。(c)说明近费米面脏极限$f$可近常数，$\Delta\to0$恢复正常电导与自旋因子。(d)化简至Mattis--Bardeen积分、数值画$\sigma_1/\sigma_N$并讨论；要求$\Delta\tau\ll1$。
\end{exercise}
\begin{solution}
(a)Bogoliubov展开电流，作用于零准粒子真空仅留下双创建项；速度矩阵在时间反演下变号，电子与空穴振幅相减，得$\sigma_1=(\pi e^2/\omega)\iint(uv'-vu')^2f(\xi,\xi')\delta(\omega-E-E')\dd\xi\dd\xi'$。自旋及双计数依原归一计入$f$和两创建通道。

(b)展开平方，使用$u^2,v^2$与$uv=\Delta/(2E)$，直接得所列结果。(c)正常态轨道的速度相关能量宽度约$1/\tau$，若$|\xi-\xi'|\ll1/\tau$，脏金属近似可将$f$固定在费米面值。$\Delta=0$时相干因子只保留$\xi,\xi'$异号的初末态，积分长度为$\omega$，恰恢复有限$\sigma_N$。四象限对称求和后$\xi\xi'$项抵消，剩题示$1-\Delta^2/(EE')$。

(d)改变量$\dd\xi=E\dd E/\sqrt{E^2-\Delta^2}$，$\delta$积分给
\begin{equation}
\frac{\sigma_1(\omega)}{\sigma_N}=\frac{\theta(\omega-2\Delta)}{\omega}\int_\Delta^{\omega-\Delta}\frac{E(\omega-E)-\Delta^2}{\sqrt{E^2-\Delta^2}\sqrt{(\omega-E)^2-\Delta^2}}\dd E.
\end{equation}
取$E=\Delta+(\omega-2\Delta)\sin^2\theta$可消端点平方根奇异，随卷脚本实际积分并画图。$\omega<2\Delta$没有零温吸收；阈值处两态密度虽发散，相干因子同时压低，整体从0起；高频在仍满足脏近似窗口内趋近1。零频超流$\delta$项与虚部响应不在这条有限频率吸收曲线中。
\end{solution}
\begin{center}\begin{tikzpicture}\begin{axis}[width=11cm,height=5.7cm,xlabel={$\omega/\Delta$},ylabel={$\sigma_1/\sigma_N$},title={零温脏极限光吸收},xmin=0,ymin=0,ymax=1.05]
\addplot[thick,color=structurecolor] coordinates {(0,0) (0.17272727,0) (0.34545455,0) (0.51818182,0) (0.69090909,0) (0.86363636,0) (1.0363636,0) (1.2090909,0) (1.3818182,0) (1.5545455,0) (1.7272727,0) (1.9,0) (2,0) (2.06,0.045415756) (2.12,0.087636604) (2.18,0.12696447) (2.24,0.16366573) (2.3,0.19797619) (2.36,0.23010529) (2.42,0.26023955) (2.48,0.28854562) (2.54,0.31517274) (2.6,0.34025498) (2.66,0.36391301) (2.72,0.38625572) (2.78,0.40738163) (2.84,0.42738) (2.9,0.44633191) (2.96,0.46431117) (3.02,0.48138505) (3.08,0.49761502) (3.14,0.51305731) (3.2,0.52776348) (3.26,0.54178084) (3.32,0.55515291) (3.38,0.56791976) (3.44,0.58011835) (3.5,0.59178281) (3.56,0.60294469) (3.62,0.61363321) (3.68,0.62387547) (3.74,0.63369662) (3.8,0.64312003) (3.86,0.65216745) (3.92,0.66085913) (3.98,0.66921395) (4.04,0.67724955) (4.1,0.68498239) (4.16,0.69242788) (4.22,0.69960045) (4.28,0.7065136) (4.34,0.71318) (4.4,0.71961157) (4.46,0.72581946) (4.52,0.73181419) (4.58,0.73760565) (4.64,0.74320316) (4.7,0.74861548) (4.76,0.75385089) (4.82,0.75891722) (4.88,0.76382183) (4.94,0.76857172) (5,0.77317346) (5.06,0.77763331) (5.12,0.78195719) (5.18,0.78615068) (5.24,0.79021912) (5.3,0.79416753) (5.36,0.79800072) (5.42,0.80172321) (5.48,0.80533934) (5.54,0.80885321) (5.6,0.81226873) (5.66,0.81558962) (5.72,0.81881943) (5.78,0.82196154) (5.84,0.82501917) (5.9,0.82799538) (5.96,0.8308931) (6.02,0.83371515) (6.08,0.83646418) (6.14,0.83914276) (6.2,0.84175333) (6.26,0.84429821) (6.32,0.84677966) (6.38,0.84919979) (6.44,0.85156067) (6.5,0.85386426) (6.56,0.85611243) (6.62,0.85830699) (6.68,0.86044967) (6.74,0.86254213) (6.8,0.86458595) (6.86,0.86658268) (6.92,0.86853376) (6.98,0.87044063) (7.04,0.87230463) (7.1,0.87412707) (7.16,0.87590919) (7.22,0.87765221) (7.28,0.87935729) (7.34,0.88102554) (7.4,0.88265805) (7.46,0.88425583) (7.52,0.88581991) (7.58,0.88735123) (7.64,0.88885072) (7.7,0.89031929) (7.76,0.89175779) (7.82,0.89316705) (7.88,0.89454789) (7.94,0.89590107) (8,0.89722734) (8.5,0.90732454) (9,0.91595542) (9.5,0.92339547) (10,0.92985773) (10.5,0.93550917) (11,0.94048223) (11.5,0.94488308) (12,0.94879772) (12.5,0.95229638) (13,0.95543696) (13.5,0.9582675) (14,0.96082819) (14.5,0.96315281) (15,0.96527) (15.5,0.96720412) (16,0.96897603) (16.5,0.97060366) (17,0.97210251) (17.5,0.97348605) (18,0.97476599) (18.5,0.97595258) (19,0.97705484) (19.5,0.9780807) (20,0.97903717)};
\end{axis}\end{tikzpicture}\end{center}

\originalsection{作业8（2009年4月13日截止）}
\begin{exercise}[原题1、2(a)--(c)]\label{p:512-8-1}
近$T_c$时Gaussian制GL自由能$F=\Delta f\int\dd r[t|\psi|^2+|\psi|^4/2+\xi^2|(-\ii\nabla+2eA/\hbar c)\psi|^2]$，$t=(T-T_c)/T_c$；原题取$\hbar=1$。(1)求超导--正常自由能差与热力学场$H_c$。(2a)变分求GL方程。(2b)从正常态降场，利用线性化Landau级求$H_{c2}$。(2c)求London深度$\lambda(T)$，以$\xi(T)=\xi/\sqrt{-t}$与$\lambda$表示$H_c$，证明$H_{c2}>H_c$等价$\kappa>1/\sqrt2$。
\end{exercise}
\begin{solution}
均匀极小$|\psi|^2=-t$给$\delta f=-\Delta f\,t^2/2$，故$H_c^2/(8\pi)=\Delta f\,t^2/2$。变分$\xi^2(-\ii\nabla+2eA/\hbar c)^2\psi+t\psi+|\psi|^2\psi=0$。最低磁动能本征值$\xi^2(2eH/\hbar c)$，当与$-t$抵消时正常态首次失稳，$H_{c2}=(-t)\Phi_0/(2\pi\xi^2)=\Phi_0/[2\pi\xi(T)^2]$。

梯度项二阶$A$系数为$\Delta f\xi^2(2e/\hbar c)^2|\psi|^2$，电流变分配Maxwell得$\lambda^{-2}=32\pi e^2\Delta f\xi^2(-t)/(\hbar^2c^2)$。因此$H_c=\Phi_0/[2\sqrt2\pi\lambda(T)\xi(T)]$，$H_{c2}/H_c=\sqrt2\lambda/\xi(T)$，得到临界$\kappa$。此比较是GL近$T_c$判据，不能直接套到温度很低且GL参数强变的任意材料。
\end{solution}
